\section*{Übersicht verwendeter Hilfsmittel}
\markboth{Übersicht verwendeter Hilfsmittel}{Übersicht verwendeter Hilfsmittel}

Im Rahmen der Erstellung dieser Arbeit wurden neben klassischen Hilfsmitteln (Literaturquellen, Schreib- und Satzprogrammen) auch generative IT-/KI-Werkzeuge unterstützend eingesetzt.

\medskip
Die verwendeten Werkzeuge im Einzelnen:

\begin{itemize}
  \item \textbf{Claude Code (\url{https://claude.ai}):}
  Unterstützung bei der Implementierung des CAP-Services, der CDS-Modelle und der Workflow-Integration. Das Werkzeug wurde zur Codegenerierung, Fehleranalyse und iterativen Verbesserung der Anwendung eingesetzt. Alle generierten Codeartefakte wurden geprüft, angepasst und im Kontext der Anwendungsarchitektur validiert. Darüber hinaus wurde Claude Code zur sprachlichen Unterstützung bei der Erstellung dieser Hausarbeit verwendet.

  \item \textbf{SAP Business Application Studio:}
  Cloudbasierte Entwicklungsumgebung für die Erstellung und das Testen der CAP-Anwendung sowie der Fiori-Elements-Oberfläche.

  \item \textbf{SAP Build Process Automation:}
  Low-Code-Werkzeug zur Modellierung des Genehmigungsworkflows mit Formularen, Bedingungen und API-Aktionen.

  \item \textbf{Overleaf / \LaTeX{}:}
  Satz und Formatierung der vorliegenden Hausarbeit.
\end{itemize}
